\theory{Queueing theory}{How many servers, machines, or communication channels are needed so that waiting remains acceptable without wasting capacity?}{Queueing theory models arrivals, service, waiting discipline, and departures as random processes. It predicts queue lengths, waiting times, utilization, throughput, and the probability of delay.}{Telephone networks, Internet routers, hospitals, factories, shops, traffic systems, call centers, cloud computing, and project management.}{Arrival and service processes determine waiting-time distributions, utilization, stability, and capacity in a stochastic system.}

\paragraph{The problem and history}
Every queue has arrivals, one or more servers, a waiting area, a service rule, and departures. The same structure describes a supermarket line, packets entering a router, patients entering an emergency department, and jobs sent to a processor. The goal is to calculate distributions and averages that support capacity decisions \cite{queue_ref1}.

Agner Krarup Erlang created the subject while studying incoming calls at the Copenhagen Telephone Exchange. He modeled arrivals by a Poisson process and solved early single- and multiple-server systems. Pollaczek, Khinchin, Kendall, and Kingman developed major later results. Kleinrock applied queueing models to message switching and packet networks \cite{queue_ref2,queue_ref3}.

\end{historylist}
\section*{3. Theory}
\subsection*{Core theory}
\paragraph{Building a model}
An arrival process describes when customers or jobs appear. A Poisson process with rate $\lambda$ has independent arrivals and exponentially distributed interarrival times. The service process describes how long a server takes; exponential service is denoted by M, deterministic service by D, and a general distribution by G. Kendall's notation writes a queue as A/S/$c$, where $c$ is the number of servers. M/M/1 has Poisson arrivals, exponential service, and one server.

The queue may have finite or infinite capacity, a finite or infinite customer population, and different disciplines. First-come first-served is common, but priority, last-come first-served, processor sharing, and random selection also occur. A customer may abandon the queue or be blocked when no buffer is free \cite{queue_ref1,queue_ref4}.

\paragraph{The M/M/1 queue}
Let arrivals occur at rate $\lambda$ and service at rate $\mu$. The traffic intensity is
\[
\rho=\frac{\lambda}{\mu}.
\]
For an unlimited queue to have a steady state, $\rho<1$ is required. If $\rho\geq1$, work arrives as fast as or faster than it can be processed.

When $\rho<1$, the probability of finding $n$ customers in the system is
\[
P(N=n)=(1-\rho)\rho^n,\qquad n=0,1,2,\ldots.
\]
The mean number in the system is $L=\rho/(1-\rho)$, and the mean time in the system is $W=1/(\mu-\lambda)$. The mean number waiting is $L_q=\rho^2/(1-\rho)$ and the mean waiting time is $W_q=\rho/(\mu-\lambda)$ \cite{queue_ref1}.

\paragraph{Little's law and networks}
Little's law states
\[
L=\lambda W,
\]
where $L$ is the long-run average number in the system, $\lambda$ is throughput, and $W$ is the average time spent in the system. If a website completes $600$ requests per minute and the average request takes $0.2$ minutes, the average number in the system is $120$.

Queueing networks connect several nodes by routing rules. One may seek an equilibrium distribution, but transient behavior matters when a hospital faces a sudden surge or a computer system starts after a failure. Product-form networks have equilibrium probabilities that factor into contributions from individual nodes \cite{queue_ref2,queue_ref3}.

\paragraph{Multiple servers and variability}
An M/M/c queue has $c$ identical servers. Adding a server usually reduces delay, but the benefit depends on utilization and variability. Deterministic service is less disruptive than highly variable service with the same mean. The Pollaczek--Khinchin formula for an M/G/1 queue makes the effect of service-time variance explicit.

Kingman's approximation for a G/G/1 queue shows that mean waiting grows with utilization, arrival variability, and service variability. This explains why a system operating at ninety-nine percent capacity can be much worse than one operating at eighty percent: random bursts have nowhere to go \cite{queue_ref4,queue_ref5}.

\paragraph{How it solves real problems}
In a hospital, arrivals vary by hour and service times vary by patient. A model estimates the probability of a long wait under one staffing plan, then compares it with plans using another doctor or a fast-treatment stream. The output is a tradeoff between cost, congestion, and service quality.

In a router, packets arrive and compete for a finite buffer. If the service rate is too low, the buffer fills and packets are dropped. Queueing analysis estimates delay and loss, helping engineers choose link speeds and buffer sizes. In cloud computing, the same calculation determines how many virtual machines are needed for a target response time \cite{queue_ref2}.

\section*{4. Important theorems and results}
\begin{enumerate}[leftmargin=*,itemsep=0.35em]
\item \textbf{Erlang formulas.} Erlang's loss and delay formulas give blocking and waiting probabilities in telephone and multi-server systems.

\item \textbf{M/M/1 stationary distribution.} When $\rho<1$, $P(N=n)=(1-\rho)\rho^n$.

\item \textbf{Little's law.} In a stable long-run system, $L=\lambda W$.

\item \textbf{Pollaczek--Khinchin formula.} The M/G/1 mean waiting time depends on the first two moments of service time, so variability increases delay.

\item \textbf{Kingman's formula.} For G/G/1 queues, mean waiting is approximately proportional to $\rho/(1-\rho)$ and to the sum of arrival and service variability.

\item \textbf{Birth--death balance.} A single queue can be represented by transitions from $n$ to $n+1$ at arrival rate and from $n$ to $n-1$ at service rate.
\end{enumerate}
\noindent\emph{Source for these results:} \cite{queue_ref1}.

\theoryapplications
\theoryconnections{queue_theory}{%
\item Probability describes random arrivals and service times, while a stochastic process records how the number of waiting jobs changes through time.
\item In a memoryless model, the current queue length is enough to predict the next transition, producing a Markov chain and balance equations.
\item Operations research then turns the resulting waiting-time and utilization formulas into a staffing or capacity decision \cite{queue_ref1,queue_ref2}.
}{%
\item Queueing theory helps a call center choose how many agents are needed to keep waiting time below a target, and helps a data center decide how much service capacity prevents overload.
\item The same calculation estimates traffic congestion, hospital beds, factory buffers, and network delay.
\item Its distinctive contribution is to show the cost of variability: even when average arrival and service rates are equal, random bursts can make the queue grow without bound \cite{queue_ref1,queue_ref3}.
}

\section*{Glossary}
\begin{description}[leftmargin=*,style=nextline,itemsep=0.3em]
\item[\textbf{Traffic intensity.}] The ratio $\rho=\lambda/\mu$ of the arrival rate to the service rate; a queue with unlimited capacity reaches a steady state only when $\rho<1$.
\item[\textbf{Poisson process.}] A model of random arrivals in which interarrival times are independent and exponentially distributed, characterized by a constant average rate $\lambda$.
\item[\textbf{Little's law.}] The relationship $L=\lambda W$ stating that the long-run average number in a stable system equals the throughput multiplied by the average time spent in the system.
\item[\textbf{M/M/1 queue.}] The simplest queueing model with Poisson arrivals, exponentially distributed service times, one server, and unlimited waiting capacity.
\item[\textbf{Kendall's notation.}] A shorthand A/S/$c$ that classifies a queue by its arrival distribution, service distribution, and number of servers.
\item[\textbf{Steady state.}] A long-run equilibrium condition in which the probability distribution of the number in the system no longer changes with time.
\item[\textbf{Pollaczek--Khinchin formula.}] An exact expression for the mean waiting time in an M/G/1 queue showing that service-time variance directly increases delay.
\item[\textbf{Kingman's formula.}] An approximation for the mean waiting time in a G/G/1 queue proportional to $\rho/(1-\rho)$ and to the sum of arrival and service variabilities.
\item[\textbf{Interarrival time.}] The time between consecutive arrivals; exponentially distributed in a Poisson process.
\item[\textbf{Service time.}] The duration a server takes to complete one job.
\item[\textbf{Utilization.}] The fraction of time a server is busy, equal to $\rho$ in single-server models.
\item[\textbf{Birth--death process.}] A continuous-time Markov chain where transitions only increase or decrease the state by one.
\item[\textbf{Erlang formula.}] Formulas giving blocking and waiting probabilities in multi-server telephone systems.
\end{description}
\section*{7. Sample Questions}
\emph{Exercise reference:} \cite{queue_ref1}. The cited edition is the source for the exercise set; consult its Exercises section for the official statement and numbering.
\begin{enumerate}[leftmargin=*,itemsep=0.9em]
\item \textbf{Exercise 1.} An M/M/1 queue has arrival rate $\lambda=4$ jobs per minute and service rate $\mu=5$ jobs per minute. Compute (a) the traffic intensity $\rho$, (b) the mean number in the system $L$, and (c) the probability that the system is empty. \emph{Solution.} (a) $\rho=\lambda/\mu=4/5=0.8$. (b) $L=\rho/(1-\rho)=0.8/0.2=4$. (c) $P(N=0)=1-\rho=0.2$.

\item \textbf{Exercise 2.} For an M/M/1 queue with $\lambda=6$ per hour and $\mu=10$ per hour, compute the probability of having exactly $3$ customers in the system in steady state. \emph{Solution.} $\rho=6/10=0.6$. The steady-state probability is $P(N=3)=(1-\rho)\rho^3=0.4\times 0.6^3=0.4\times 0.216=0.0864$.

\item \textbf{Exercise 3.} A call center completes $200$ calls per hour, and the average call duration is $4$ minutes. Using Little's law, find the average number of calls in progress simultaneously. \emph{Solution.} Convert to consistent units: $\lambda=200/60=10/3$ calls per minute, $W=4$ minutes. Little's law gives $L=\lambda W=(10/3)\times 4=40/3\approx 13.3$ calls.

\item \textbf{Exercise 4.} In an M/M/2 queue with $\lambda=8$ per minute and $\mu=5$ per minute per server, compute the traffic intensity per server and determine whether the queue is stable. \emph{Solution.} The offered load is $a=\lambda/\mu=8/5=1.6$. With $c=2$ servers, the per-server utilization is $\rho=a/c=1.6/2=0.8<1$, so the queue is stable.

\item \textbf{Exercise 5.} For an M/M/1 queue with $\lambda=3$ per minute and $\mu=4$ per minute, compute the mean waiting time in queue (not including service). \emph{Solution.} $\rho=\lambda/\mu=3/4=0.75$. The mean number waiting is $L_q=\rho^2/(1-\rho)=0.75^2/0.25=0.5625/0.25=2.25$. By Little's law applied to the queue, $W_q=L_q/\lambda=2.25/3=0.75$ minutes, or $45$ seconds. Equivalently, $W_q=\rho/(\mu-\lambda)=0.75/1=0.75$ minutes.

\end{enumerate}
\section*{8. References}
\begin{thebibliography}{9}
\bibitem{queue_ref1} D. Gross, J. F. Shortle, J. M. Thompson, and C. M. Harris, \emph{Fundamentals of Queueing Theory}, 4th ed., Wiley, 2008.
\bibitem{queue_ref2} L. Kleinrock, \emph{Queueing Systems, Volume 1: Theory}, Wiley, 1975.
\bibitem{queue_ref3} L. Kleinrock, \emph{Queueing Systems, Volume 2: Computer Applications}, Wiley, 1976.
\bibitem{queue_ref4} J. W. Cohen, \emph{The Single Server Queue}, North-Holland, 1969.
\bibitem{queue_ref5} J. F. C. Kingman, \emph{Heavy Traffic in Single Server Queues}, Notes in Applied Mathematics, American Mathematical Society, 1965.
\end{thebibliography}
